\hypersetup{pageanchor=false}
\begin{titlepage}
\color{ink}
\vspace*{12mm}
{\large\sffamily LEARNING NOTES \quad / \quad QKSITFALL}\par
\vspace{16mm}
{\Huge\bfseries Quantum EV Charging}\par
\vspace{4mm}
{\LARGE From booking rules to QAOA}\par
\vspace{10mm}
\rule{\textwidth}{0.8pt}
\vspace{5mm}
{\large A software engineer's guide to the actual Python implementation}\par
\vspace{10mm}
This document explains the complete project README, including the mathematical
steps hidden inside the implementation. It begins with ordinary booking data,
builds the optimization objective, translates it into quantum operations, and
explains what the measurements mean.

No previous quantum mechanics is assumed. Symbols are introduced before use;
longer derivations, a glossary, and exercises with answers appear in the appendices.
\vfill
\textbf{Execution model}\par
Qiskit QAOA executed on a local Aer simulator. Classical Python prepares the
problem and adjusts circuit angles. No IBM API key is used by the current command.
\vspace{8mm}
\textbf{Code snapshot}\par
1 November 2026. Qiskit 2.6.1; Qiskit Aer 0.18.2; SciPy 1.18.1.
The numerical examples follow \code{main.py} and \code{src/quantum_scheduler.py}.
\end{titlepage}
\hypersetup{pageanchor=true}
\pagenumbering{roman}

\tableofcontents
\clearpage
\pagenumbering{arabic}

\section{How to use this guide}
Read Sections~\ref{sec:physical}--\ref{sec:qubo} first if the scheduling model or
penalty expressions are unclear. Read Sections~\ref{sec:vocabulary}--\ref{sec:gates}
slowly for the transition from bits to gates. Then run the program while reading
Sections~\ref{sec:training}--\ref{sec:experiments}.

The implemented workflow is
\begin{equation*}
\begin{gathered}
\text{vehicle, station, and travel data}
\longrightarrow \text{possible bookings}
\longrightarrow \text{binary objective}\\
\longrightarrow \text{quantum cost operator}
\longrightarrow \text{QAOA circuit}
\longrightarrow \text{measured booking bits}.
\end{gathered}
\end{equation*}
The quantum circuit participates in searching for a schedule. It does not
calculate road distances, parse vehicle data, or invent the scheduling rules.

Run these commands from the repository root:
\begin{lstlisting}
.venv/bin/python main.py
.venv/bin/python main.py --layers 1 --show-circuit
.venv/bin/python -m unittest discover -s tests -v
\end{lstlisting}
The main entry point prints six quantum checkpoints. The document follows those
checkpoints, with additional foundations where they are needed.

\begin{keyidea}
There are three different objects to keep separate: a real booking, the binary
variable that selects it, and the qubit used to represent that variable.
\end{keyidea}

\section{The charging problem before quantum computing}\label{sec:physical}
\subsection{The input scenario}
We coordinate two cars and two stations. Each station has one charging port
with power $11\,\mathrm{kW}$. Each reservation occupies an entire one-hour slot:
08:00--09:00 or 09:00--10:00.

\begin{center}
\begin{tabular}{lrr}
\toprule
Vehicle property & Car A & Car B\\
\midrule
Battery capacity & $50\,\mathrm{kWh}$ & $50\,\mathrm{kWh}$\\
Initial charge & $30\%$ & $12\%$\\
Target charge & $50\%$ & $30\%$\\
Consumption & $0.2\,\mathrm{kWh/km}$ & $0.2\,\mathrm{kWh/km}$\\
Available to depart & 07:45 & 07:50\\
Charging deadline & 10:00 & 09:00\\
\bottomrule
\end{tabular}
\end{center}

\begin{center}
\begin{tabular}{lrrrr}
\toprule
Journey & A--North & A--South & B--North & B--South\\
\midrule
Distance (km) & 2 & 5 & 1 & 20\\
Duration (min) & 10 & 15 & 5 & 30\\
\bottomrule
\end{tabular}
\end{center}

Electricity at North costs Rs~2/kWh in the early slot and Rs~1/kWh in the late
slot. South costs Rs~3/kWh and Rs~2/kWh respectively. ``Rs'' denotes rupees;
we keep monetary amounts separate from dimensionless optimization score points.

\subsection{Energy, power, and arrival}
Let $B_v$ be vehicle $v$'s battery capacity in kWh; $s_v$ and $t_v$ its initial
and target state-of-charge \emph{fractions}; $d_{vs}$ its distance to station $s$;
and $e_v$ its consumption in kWh/km. The Python inputs store percentages, so
$30\%$ is entered as \code{30} and divided by $100$ in the calculation.
\begin{align}
E_v^{\mathrm{initial}} &= B_v s_v,\\
E_{vs}^{\mathrm{drive}} &= d_{vs}e_v,\\
E_{vs}^{\mathrm{arrival}} &= E_v^{\mathrm{initial}}-E_{vs}^{\mathrm{drive}},\\
D_{vs} &= B_v t_v-E_{vs}^{\mathrm{arrival}}.
\end{align}
Here $D_{vs}$ is the charging energy required \emph{at that station}. A farther
station can require more charging, because the drive reduces the arrival energy.

For A driving to North:
\begin{align*}
E_A^{\mathrm{initial}} &= 50(0.30)=15.0\,\mathrm{kWh},\\
E_{A,N}^{\mathrm{drive}} &= 2(0.2)=0.4\,\mathrm{kWh},\\
E_{A,N}^{\mathrm{arrival}} &=15.0-0.4=14.6\,\mathrm{kWh},\\
D_{A,N} &=50(0.50)-14.6=10.4\,\mathrm{kWh}.
\end{align*}
If station power is $W_s$ kW, ideal charging duration is
\begin{equation}
\tau_{vs}=60\frac{D_{vs}}{W_s}\quad\text{minutes}.
\end{equation}
Thus A needs $60(10.4/11)\approx56.73$ minutes at North. Energy in kWh is an
amount; power in kW is an amount delivered per hour.

\subsection{Individually feasible bookings}
The function \code{generate_options()} examines every vehicle--station--slot
combination. A booking survives if all four checks pass:
\begin{enumerate}
\item Arrival battery is at least the reserve, $0.05B_v$ by default.
\item Arrival time is no later than the selected slot's start.
\item Starting then, actual charging completes no later than the vehicle deadline.
\item The required energy fits within the slot's charging capacity.
\end{enumerate}
For B at South, arrival energy would be $6-20(0.2)=2\,\mathrm{kWh}$, below the
$2.5\,\mathrm{kWh}$ reserve. B also cannot take North's late slot because charging
would finish after its 09:00 deadline. Only North's early slot survives for B.

For booking $i$, define electricity bill $b_i$, waiting time $w_i$ in minutes,
and booking score $c_i$:
\begin{equation}\label{eq:booking-cost}
b_i=D_{vs}\,r_{s,k},\qquad
w_i=T_k^{\mathrm{start}}-T_{vs}^{\mathrm{arrival}},\qquad
c_i=b_i+\lambda w_i,
\end{equation}
where $r_{s,k}$ is the electricity tariff and $\lambda=0.25$ is the waiting weight.
Numerically, each rupee contributes one score point and each waiting minute
contributes $0.25$ points. The score is a preference, not a physical law.

\section{Checkpoint 1: turn possible bookings into bits}
The five surviving options are assigned consecutive indices by
\code{build_qubo()}. Each $x_i\in\{0,1\}$ selects one booking. Its corresponding
qubit is $q_i$.
\begin{center}
\begin{tabular}{clrrr}
\toprule
Bit & Booking & Bill (Rs) & Wait (min) & $c_i$\\
\midrule
$x_0$ & A, North, 08:00 & 20.80 & 5 & 22.05\\
$x_1$ & A, North, 09:00 & 10.40 & 65 & 26.65\\
$x_2$ & A, South, 08:00 & 33.00 & 0 & 33.00\\
$x_3$ & A, South, 09:00 & 22.00 & 60 & 37.00\\
$x_4$ & B, North, 08:00 & 18.40 & 5 & 19.65\\
\bottomrule
\end{tabular}
\end{center}
All five variables appear in the objective, but a valid schedule selects only
one booking for each car. The vector
\begin{equation}
\boldsymbol{x}=(0,1,0,0,1)
\end{equation}
selects A at North late and B at North early. Their reservations occur at
different times. We plan both cars jointly because their options share resources.

The groups are built from data:
\begin{lstlisting}[language=Python]
groups = {"A": [0, 1, 2, 3], "B": [4]}
\end{lstlisting}
This displayed dictionary is the sample result, not a hardcoded assumption in
the builder. Adding cars or options changes the groups and the number of qubits.
B's forced booking could be fixed classically, but the current example retains
its qubit to make the rule application uniform.

\section{Checkpoint 2: build the QUBO objective}\label{sec:qubo}
QUBO stands for \emph{quadratic unconstrained binary optimization}. ``Binary''
means bits, ``quadratic'' means terms involving at most pairs of bits, and
``unconstrained'' means the scheduling rules have been incorporated as penalties
inside the objective. It does not mean the scheduling problem has no rules.

\subsection{Two reusable scheduling rules}
Let $\mathcal V$ be the vehicle set, $G_v$ the indices of bookings belonging to
vehicle $v$, and $\mathcal C$ the set of unordered pairs of bookings for different
vehicles that occupy the same station and slot. The general objective is
\begin{equation}\label{eq:general-qubo}
f(\boldsymbol{x})=
\underbrace{\sum_i c_i x_i}_{\text{booking scores}}
+\underbrace{P\sum_{v\in\mathcal V}\left(\sum_{i\in G_v}x_i-1\right)^2}_{\text{one booking per car}}
+\underbrace{P\sum_{(i,j)\in\mathcal C}x_i x_j}_{\text{no shared charger conflicts}}.
\end{equation}
The sum over a car counts its selected bookings. Its penalty is zero exactly
when that count equals one. A conflict product is one only if both conflicting
bookings are selected. These templates apply to any number of cars and options
under the current single-port, single-slot-session model.

\begin{center}
\begin{tabular}{crrrr}
\toprule
Selected bookings for one car & 0 & 1 & 2 & 3\\
\midrule
Exactly-one penalty & $P$ & $0$ & $P$ & $4P$\\
\bottomrule
\end{tabular}
\end{center}

\subsection{The actual sample equation}
Here $G_A=\{0,1,2,3\}$, $G_B=\{4\}$, and $\mathcal C=\{(0,4)\}$. Therefore
\begin{align}\label{eq:sample-qubo}
f(\boldsymbol{x})={}&22.05x_0+26.65x_1+33x_2+37x_3+19.65x_4\notag\\
&+P(x_0+x_1+x_2+x_3-1)^2\notag\\
&+P(x_4-1)^2+Px_0x_4.
\end{align}
The first penalty belongs to A; the second belongs to B; the last belongs to
their shared North/early resource. Giving B more options changes its group sum.
Another conflict creates another pair term automatically.

For the valid vector $(0,1,0,0,1)$ every penalty is zero:
\begin{equation}
f(0,1,0,0,1)=26.65+19.65=46.30.
\end{equation}
For $(1,0,0,0,1)$ both cars receive a booking, but both use North early:
\begin{equation}
f(1,0,0,0,1)=22.05+19.65+P=41.70+P.
\end{equation}
The per-car penalties alone would not detect this charger conflict.

\subsection{Choosing the penalty without solving the schedule}
All booking scores are nonnegative. Define the upper bound
\begin{equation}\label{eq:penalty-bound}
U=\sum_{v\in\mathcal V}\max_{i\in G_v}c_i,
\qquad P=U+1.
\end{equation}
Any valid schedule selects one option per car and therefore scores at most $U$.
Any invalid bit string incurs at least $P$ in penalties plus nonnegative booking
cost. If a valid schedule exists, no invalid string can be the global minimizer.
This statement is about the objective, not a promise that QAOA finds its minimum.

For the sample, $U=37+19.65=56.65$ and $P=57.65$. A custom smaller penalty may
make invalid strings attractive. Increasing $P$ far beyond necessity can make
the differences between valid operating costs harder to resolve during training.

\subsection{How the builder stores the equation}
Using $x_i^2=x_i$ for bits, the builder expands every group penalty:
\begin{equation}\label{eq:group-expansion}
P\left(\sum_{i\in G_v}x_i-1\right)^2
=P-P\sum_{i\in G_v}x_i
+2P\sum_{\substack{i<j\\i,j\in G_v}}x_i x_j.
\end{equation}
The stored representation is
\begin{equation}\label{eq:expanded-qubo}
f(\boldsymbol{x})=C_0+\sum_i L_i x_i+\sum_{i<j}Q_{ij}x_i x_j.
\end{equation}
The code names are \code{constant}, \code{linear}, and \code{quadratic}.
Here $C_0=P|\mathcal V|$, each $L_i=c_i-P$, each same-car pair gets $2P$, and
each different-car resource conflict gets $P$. Negative $L_i$ values are a
consequence of expanding penalties; they are not negative electricity bills.

\begin{checkpoint}
If a third car has bookings $x_5$ and $x_6$, write its exactly-one penalty.
If $x_1$ and $x_5$ share a charger and slot, write their conflict penalty.
Answers appear in Appendix~\ref{app:answers}.
\end{checkpoint}
